\title{Parameter-Efficient Orthogonal Finetuning via Factorization}

\begin{abstract}
Large foundational models have become widespread, yet training them from the ground up is prohibitively costly. Consequently, the efficient adaptation of these robust models to specific downstream tasks has gained importance. In this work, we explore a systematic finetuning approach — Orthogonal Finetuning (OFT) — for adapting to downstream tasks. Although OFT has shown strong generalization capabilities, it still involves a considerable number of trainable parameters due to the high dimensionality of orthogonal matrices. To tackle this, we begin by analyzing OFT through the lens of information transmission and identify several essential requirements to improve parameter efficiency. Drawing inspiration from the Cooley-Tukey fast Fourier transform algorithm, which facilitates efficient information flow, we propose a compact orthogonal parameterization employing butterfly structures. We integrate this parameterization into OFT, resulting in a new parameter-efficient finetuning technique termed Orthogonal Butterfly (BOFT). Since BOFT subsumes OFT as a particular instance, it establishes a broader framework for orthogonal finetuning. Finally, we perform a thorough empirical evaluation of adapting large vision transformers, large language models, and text-to-image diffusion models across a variety of vision and language downstream tasks.
\end{abstract}

\section{Introduction}

Recent developments such as ChatGPT~\cite{brown2020language,chen2021evaluating} and Stable Diffusion~\cite{rombach2022high} showcase the impressive generalization capabilities of large foundational models. The rapid improvements in these models’ performance have been accompanied by a substantial increase in parameter counts (e.g., GPT-3 contains roughly 175B parameters). As a consequence, training such foundation models from scratch has become increasingly impractical for most researchers. Progress in the field hence relies on the ability to adapt these models effectively without complete retraining. In other words, there is a critical need to efficiently tailor existing foundation models to downstream tasks. This can primarily be achieved through three approaches: \emph{model finetuning}~\cite{hulora2022,zaken2022bitfit,guo2020parameter,raffel2020exploring,qiu2023controlling,zhang2022adaptive,chavan2023one}, where the model structure remains fixed and only a subset of parameters is fine-tuned; \emph{adapter tuning}~\cite{rebuffi2017learning,houlsby2019parameter,pfeiffer2020adapterfusion,lin2020exploring,he2021towards}, which involves adding new trainable parameters to the original model; and \emph{prompt tuning}~\cite{li2021prefix,lester2021power}, where trainable prefix tokens are prepended to the input. Among these, model finetuning stands out as a straightforward yet effective strategy that also avoids any increase in inference latency.

The core idea behind model finetuning is to maintain similarity between the pretrained model and the finetuned model according to specific criteria, thereby preserving the knowledge acquired during pretraining. For example, many existing finetuning techniques utilize a small learning rate, which heuristically helps keep the pretrained and finetuned models closely aligned. Considering the structured nature of weight matrices, a more principled method aims to retain the relational structure within these matrices, specifically the pairwise angles between neurons. This concept motivates a new finetuning paradigm called Orthogonal Finetuning~(OFT)~\cite{qiu2023controlling}, where neurons in a given layer undergo transformation by a shared orthogonal matrix, ensuring that the pairwise angles among neurons are theoretically preserved throughout the finetuning. Although OFT has shown promising results in terms of generalization and convergence when applied to finetuning text-to-image diffusion models~\cite{qiu2023controlling}, the dimensionality of orthogonal matrices leads to a large number of trainable parameters. To mitigate this, OFT imposes a block-diagonal structure on the orthogonal matrix, reducing parameter count. Nevertheless, this parameter efficiency incurs drawbacks: the orthogonal matrix adopts a fixed sparsity pattern and the orthogonal transformation is performed independently within each block. This block-diagonal sparsity pattern, while saving parameters, may introduce undesired inductive biases (e.g., limiting expressiveness and preventing approximation of standard linear transformations). Crucially, determining an optimal sparsity pattern remains an open challenge.

The crucial challenge in solving this problem lies in generating a dense orthogonal matrix while maintaining parameter efficiency. Although this may initially appear impractical since a dense orthogonal matrix of dimension $d$ typically requires $\mathcal{O}(d^2)$ parameters, we adopt an innovative strategy by constructing a dense orthogonal matrix through the product of several factorized sparse matrices. This method is motivated by the observation that the number of trainable parameters can be decreased by exchanging memory usage for additional computation. Because the orthogonal matrix is represented as a product of sparse matrices, these multiplications must be executed repeatedly during each training step. To better understand the matrix factorization, we frame the creation of a dense orthogonal matrix as an information transmission problem. From this viewpoint, generating a dense orthogonal matrix via multiplication of sparse matrices can be interpreted as transmitting information over a grid-structured graph. This framework for information transmission allows us to devise various sparse matrix factorizations that minimize trainable parameters while retaining enough expressive power to form dense matrices.

To ensure parameter efficiency within this information transmission framework, we take inspiration from the butterfly structures used in the Cooley-Tukey fast Fourier transform algorithm~\cite{cooley1965algorithm}, where data exchange among different nodes is highly efficient~\cite{klasing1994broadcasting}. Specifically, the butterfly graph structure intrinsic to the Cooley-Tukey algorithm naturally induces a sparse matrix factorization technique known as butterfly factorization. Utilizing butterfly factorization, we can create a $d$-dimensional dense matrix as the product of $\mathcal{O}(\log d)$ sparse matrices, each containing only $\mathcal{O}(d)$ non-zero elements. By guaranteeing that every sparse matrix is orthogonal, we derive an efficient orthogonal parameterization requiring just $\mathcal{O}(d \log d)$ parameters, a substantial reduction compared to the conventional full parameterization. Building upon this efficient orthogonal parameterization, we introduce a novel parameter-efficient finetuning approach called Orthogonal Butterfly (BOFT). BOFT encompasses OFT as a special case and establishes a broad orthogonal finetuning framework. A common feature of both the block-diagonal and butterfly structures is sparsity; each incorporates data sparsity into orthogonal matrices to diminish the number of trainable parameters. It is also insightful to compare our method with the low-rank structure in LoRA~\cite{hulora2022}; both low-rank and sparse matrices belong to prominent categories of structured matrices~\cite{candes2011robust} that enable parameter efficiency.

In contrast to the block-diagonal structure employed by OFT to balance expressivity and regularity, BOFT utilizes the butterfly structure to create a more gradual interpolation between matrices drawn from the full orthogonal group (expressivity) and identity matrices (regularity). This approach allows us to identify a more compact hypothesis class within the orthogonal group tailored for downstream tasks. Given the extensive use of butterfly structures in many rapid linear transformations, such as the discrete Fourier and discrete cosine transforms, we contend that our structured strategy for parameter efficiency introduces an implicit inductive bias that enhances generalizability and mitigates overfitting. Our reasoning stems from the fact that the butterfly structure can readily reconstruct numerous classical linear transforms, whereas the block-diagonal structure in OFT cannot recover any. Our key contributions are as follows:

\begin{itemize}[leftmargin=*,nosep]
    \item We address the challenge of parameter efficiency in orthogonal finetuning by introducing a novel information transmission framework. This framework recasts the design of parameter-efficient dense orthogonal matrices as an information transmission problem within a grid-structured graph.
    \item Drawing inspiration from the butterfly structures found in the Cooley-Tukey algorithm, we propose Orthogonal Butterfly, a parameter-efficient orthogonal finetuning technique, developed within the information transmission framework.
    \item We offer several theoretical insights explaining why BOFT can drastically reduce the number of trainable parameters while maintaining strong expressivity and stable training. Moreover, due to its matrix factorization, BOFT features an interesting weight interpolation property (see Figure~\ref{fig:boft_interpolation}).
    \item For the first time, we extend orthogonal finetuning beyond controllable text-to-image generation~\cite{qiu2023controlling}, demonstrating its promise as a versatile model finetuning method. To support this, we apply BOFT to various adaptation tasks spanning computer vision and natural language processing. BOFT significantly outperforms existing state-of-the-art techniques, confirming its superior parameter efficiency and generalization capability.
\end{itemize}

\section{Related Work}

\textbf{Parameter-efficient finetuning (PEFT)}. As foundational models (\emph{e.g.}, \cite{brown2020language,radford2021learning,rombach2022high,kirillov2023segment}) grow larger and more capable, there has been significant interest in finetuning these models for downstream applications in a parameter-efficient manner~\cite{treviso2023efficient,ding2023parameter,lialin2023scaling}. Among the various PEFT approaches~\cite{houlsby2019parameter,aghajanyan2020intrinsic,hulora2022,edalati2022krona,wang2022adamix,gheini2021cross,zaken2022bitfit,guo2020parameter,sung2021training,ansell2022composable,lester2021power,li2021prefix,vu2022spot,he2021towards,mao2021unipelt,karimi2021compacter,liu2022few,sung2022lst,chen2022parameter,jia2022visual,chen2022adaptformer,zhang2022neural,jie2023fact,lian2022scaling,luo2023towards}, methods based on reparameterization~\cite{aghajanyan2020intrinsic,hulora2022,edalati2022krona} are most closely related to our work. LoRA~\cite{hulora2022} finetunes the pretrained weight matrix by adding the product of two low-rank matrices, yielding strong results on natural language tasks. Since LoRA fixes the rank of all added matrices, several works \cite{zhang2022adaptive,valipour2022dylora,zhang2023increlora} dynamically adjust the rank per layer to better distribute the parameter budget. This straightforward low-rank weight reparameterization has thus gained wide adoption \cite{dettmers2023qlora,zi2023delta,chavan2023one}. Motivated by the role of hyperspherical energy in characterizing generalization~\cite{liu2018learning,liu2021orthogonal}, \cite{qiu2023controlling} introduces orthogonal finetuning (OFT) as an alternative and effective method for finetuning text-to-image diffusion models. Specifically, OFT learns an orthogonal matrix that transforms neurons within the same layer, resulting in improved generalization and consistently more stable training compared to LoRA. Despite its strong performance, OFT typically requires more trainable parameters than LoRA. Hence, enhancing the parameter efficiency of OFT is an important goal. Furthermore, it remains unknown whether OFT can be effectively applied to a broader range of adaptation tasks beyond controlling text-to-image diffusion models~\cite{qiu2023controlling}. BOFT addresses this by improving OFT’s parameter efficiency through butterfly factorization, enabling us to demonstrate the strength of orthogonal finetuning in general adaptation scenarios.

\textbf{Butterfly structures}. The radix-2 Cooley-Tukey algorithm breaks down the $N$-point discrete Fourier transform into two $\frac{N}{2}$-point transforms recursively, resulting in a butterfly structure that can be expressed as a product of several sparse matrices (often called a butterfly matrix). Butterfly matrices have been utilized to parameterize orthogonal matrices for avoiding pivoting in Gaussian elimination and boosting efficiency~\cite{parker1995random}, to stabilize training of recurrent neural networks~\cite{jing2017tunable}, and in kernel approximation~\cite{munkhoeva2018quadrature}. Works such as \cite{dao2019learning,dao2019kaleidoscope} have explored learning fast linear transforms parameterized by butterfly matrices, while \cite{chen2022pixelated,dao2022monarch} apply butterfly matrices to facilitate efficient neural network training. Butterfly structures have also proven useful in fast matrix-vector multiplication~\cite{michielssen1996multilevel,o2010algorithm}, data-sparse matrix approximation~\cite{li2015butterfly}, and network transmission~\cite{klasing1994broadcasting,li2003linear,rajasingh2016transmission}. Unlike prior work, our focus is on leveraging butterfly structures to improve the parameter efficiency of OFT.

\section{An Information Transmission View on Orthogonal Finetuning}

We begin with an overview of OFT fundamentals. To fine-tune the pretrained weight matrix, OFT redefines the new weight matrix as the product of a learnable orthogonal matrix and the fixed pretrained weight matrix. Unlike LoRA, which updates weights by adding a low-rank matrix, OFT updates weights via multiplication by an orthogonal matrix. To maintain parameter efficiency, LoRA employs a low-rank structure, whereas the original OFT utilizes a block-diagonal orthogonal structure~\cite{qiu2023controlling} (the smaller the diagonal block sizes, the higher the parameter efficiency of OFT). An intuitive comparison is illustrated in Figure~\ref{fig:comp_lora}. The rationale for applying an orthogonal transformation during weight matrix fine-tuning is to preserve the pairwise angles between neurons~\cite{liu2018learning,liu2021orthogonal,qiu2023controlling}, thereby largely retaining the semantic knowledge acquired during pretraining. Specifically, OFT optimizes an orthogonal matrix $\bm{R}\in\mathbb{R}^{d\times d}$ for a pretrained linear layer $\bm{W}^0\in\mathbb{R}^{d\times n}$ and modifies the forward computation from $\bm{z}=(\bm{W}^0)^\top \bm{x}$ to $\bm{z}=(\bm{R}\bm{W}^0)^\top \bm{x}$, where $\bm{x}\in\mathbb{R}^d$ and $\bm{z}\in\mathbb{R}^n$ are the input and output vectors, respectively. To enable zero initialization, OFT sets $\bm{R}$ as the identity matrix at the start. To guarantee that $\bm{R}$ remains orthogonal throughout fine-tuning, we adopt the Cayley parameterization as in \cite{qiu2023controlling,liu2021orthogonal}, \emph{i.e.}, $\bm{R}=(\bm{I}+\bm{Q})(\bm{I}-\bm{Q})^{-1}$, where $\bm{Q}$ is a skew-symmetric matrix satisfying $\bm{Q}=-\bm{Q}^\top$. For parameter efficiency, the block-diagonal form is introduced by expressing the orthogonal matrix $\bm{R}$ as $\text{diag}(\bm{R}_1, \bm{R}_2, \ldots, \bm{R}_r)$, where each $\bm{R}_i\in\mathbb{R}^{b\times b}, \forall i$ is a small orthogonal matrix, and $br=d$. The parameter efficiency gained via this block-diagonal approach comes at a cost—it assumes that the dimensions of a neuron (i.e., a column vector of $\bm{W}^0$) are divided into $r$ groups, and that the dimensions within each group are transformed separately with distinct orthogonal matrices. Although the block-diagonal structure has been empirically effective, partitioning a neuron's dimensions into $r$ groups based solely on their indices lacks rationale, motivating the appeal for dense orthogonal matrices. This raises a natural question: \emph{Is it possible to construct a dense orthogonal matrix while retaining parameter efficiency?}

To tackle this issue, we suggest constructing a dense orthogonal matrix by multiplying several sparse orthogonal matrices. To offer a coherent and intuitive framework for examining the sparsity structure in orthogonal matrix factorization, we reinterpret the challenge of producing a dense orthogonal matrix in OFT as an information transmission task. More precisely, generating a dense matrix $\bm{R}\in\mathbb{R}^{d\times d}$ as the product of $m$ square matrices $\thickmuskip=2mu \medmuskip=2mu \bm{R}=\bm{B}_m\bm{B}_{m-1}\cdots\bm{B}_1$ can be understood as conveying information across a grid containing $d\times (m+1)$ nodes, as shown in Figure~\ref{fig:info_trans}. This perspective arises from viewing a $d$-dimensional dense square matrix as a full connectivity map linking $d$ nodes to another set of $d$ nodes. For matrix $\bm{R}$, if the entry $R_{ij}$ equals zero, it implies that information originating from the $j$-th node cannot reach the $i$-th node. Conversely, a non-zero $R_{ij}$ signifies possible information transfer. Hence, representing the dense matrix $\bm{R}$ via the product $\bm{B}_m\bm{B}_{m-1}\cdots\bm{B}_1$ can be interpreted as sequentially exchanging information according to the graphs defined by each $\bm{B}_i, \forall i$. Information first flows through $\bm{B}_1$ and ultimately through $\bm{B}_m$. As a concrete illustration in Figure~\ref{fig:info_trans}, we examine the factorization $\bm{R}=\bm{B}_5\bm{B}_4\bm{B}_3\bm{B}_2\bm{B}_1$, with its sparsity patterns and associated graph depicted. This graph results from unfolding the matrix multiplication process. Within the induced graph, each matrix $\bm{B}_i$ represents connectivity from nodes at level $i$ to those at level $i+1$. Specifically, the $(j_1,j_2)$ element of $\bm{B}_i$ indicates whether there is a directed edge from the $j_2$-th node at level $i$ to the $j_1$-th node at level $i+1$ (with zero indicating no edge). For the product $\bm{B}_5\bm{B}_4\bm{B}_3\bm{B}_2\bm{B}_1$ to be dense, it is necessary that every node at the first level can transmit information to all nodes at the sixth level. If we focus only on $\bm{R}=\bm{B}_4\bm{B}_3\bm{B}_2\bm{B}_1$, connecting source nodes at level one and receiver nodes at level five, we observe that information from node $1$ cannot reach node $3$. Consequently, the sparsity pattern of $\bm{B}_4\bm{B}_3\bm{B}_2\bm{B}_1$ exhibits a zero at position $(3,1)$. Similar relationships between the induced graph and sparsity pattern hold for $\bm{R}=\bm{B}_3\bm{B}_2\bm{B}_1$ and $\bm{R}=\bm{B}_2\bm{B}_1$. Generally, to ensure a matrix $\bm{R}\in\mathbb{R}^{d\times d}$ is dense, the $m$ factor matrices $\bm{B}_m,\cdots,\bm{B}_1$ must correspond to a set of directed edges in a $d\times (m+1)$ grid where each directed edge connects only nodes in adjacent levels (i.e., adjacent columns), so that information from every node at the first level can be conveyed to every node at the $(m+1)$-th level.

The information transmission perspective enables us to better comprehend the sparsity pattern of factorization matrices in OFT. Figure~\ref{fig:blk_diag} illustrates the block-diagonal configuration of $\bm{R}$ in the original OFT. Although this structure reduces the number of trainable parameters, it cannot form a dense matrix $\bm{R}$. Our objective is to construct a dense orthogonal matrix by combining $m$ sparse orthogonal matrices, while minimizing the number of effective trainable parameters. From the information transmission standpoint, the general requirements for our goal are \emph{(i)} \textbf{dense connectivity}: each node at the initial level must have at least one path to every node at the final level, and \emph{(ii)} \textbf{minimum free edges}: the total edge count should be as low as possible, subject to the orthogonality constraint. It is important to note that orthogonality imposes a subtle restriction on edges between consecutive levels. For instance, for each matrix $\bm{B}_i$ to be full-rank (a prerequisite for orthogonality), there must be $d$ edges establishing a bijection between all nodes at level $i$ and all nodes at level $(i+1)$, making the minimum number of edges between adjacent levels $d$ (\emph{e.g.}, 4 in the case shown in Figure~\ref{fig:info_trans}). These $d$ edges are essential for orthogonality and are excluded from the edge count since these elements are not trainable (\emph{e.g.}, in a $d \times d$ orthogonal matrix with $d$ non-zero elements, these elements are restricted to $\pm 1$). Because orthogonal matrices require fewer parameters than full matrices, the orthogonality constraint induces additional dependencies among edges. For example, in a $2 \times 2$ orthogonal matrix, having a zero at position $(1,1)$ necessitates a zero at position $(2,2)$ as well (\emph{i.e.}, the absence of one edge can imply the absence of another). Consequently, for any feasible edge connection set, orthogonality can sometimes add or remove edges. Visualizing the non-zero pattern of the composed orthogonal matrix highlights the usefulness of the information transmission framework in OFT, as our primary concern is the presence of non-zero trainable elements in $\bm{R}$, while their exact values are irrelevant.

A straightforward dense connection between two levels involves $\mathcal{O}(d^2)$ edges (i.e., a single dense orthogonal matrix), resulting in $d^2-d$ edges (for $d=4$, this amounts to 12 edges). Figure~\ref{fig:info_trans} illustrates an example of a viable matrix factorization, which requires only $10$ edges in total, fewer than a single dense orthogonal matrix. This approach allows us to analyze the parameter efficiency of OFT through a graphical lens, making it simple to identify feasible factorizations. We draw inspiration from an intriguing topology derived from the Cooley-Tukey algorithm, known as butterfly graphs~\cite{cooley1965algorithm}, which can effectively connect $d$ source nodes to $d$ receiver nodes with just $\mathcal{O}(d\log d)$ edges. For instance, the topology shown in Figure~\ref{fig:info_trans} uses $10$ edges to achieve dense connectivity, while the butterfly network accomplishes this with only $8$ edges. In the following, we explain how the butterfly structure enhances parameter efficiency.

\section{Orthogonal Parameterization by Butterfly Factorization}

The butterfly structure was originally introduced in the Cooley-Tukey algorithm for efficiently performing the fast Fourier transform. In the Fourier transform, a local modification in the frequency domain induces a global change in the spatial domain, which conceptually aligns with our information transmission problem—each node at the first level can convey information to all nodes at the final level. Additionally, the butterfly structure has become a widely adopted computer network topology~\cite{solihin2015fundamentals,li2011linear} that facilitates efficient information exchange. Let $k \geq 2$ be a power of $2$, and define the \emph{butterfly factor} $\bm{B}^F(k)$ as

\begin{equation}\label{eq:but_factor}
\begin{aligned}
    \bm{B}^F(k) = \begin{bmatrix}
    \text{diag}(\bm{d}_1) & \text{diag}(\bm{d}_2) \\
    \text{diag}(\bm{d}_3) & \text{diag}(\bm{d}_4)
    \end{bmatrix} \in \mathbb{R}^{k \times k},
\end{aligned}
\end{equation}

where $\bm{d}_i \in \mathbb{R}^{\frac{k}{2}}, \forall i$ are vectors. For $d = 2^N$, the $d$-dimensional \emph{butterfly matrix} $\bm{B}(d) \in \mathbb{R}^{d \times d}$ is recursively defined as

\begin{equation}
\begin{aligned}
    \bm{B}(d) = \tilde{\bm{B}}(d,d) \cdot \begin{bmatrix}
    \bm{B}_1\left(\frac{d}{2}\right) & \bm{0} \\
    \bm{0} & \bm{B}_2\left(\frac{d}{2}\right)
    \end{bmatrix} = \tilde{\bm{B}}(d,d) \tilde{\bm{B}}(d, \tfrac{d}{2}) \cdots \tilde{\bm{B}}(d, 2),
\end{aligned}
\end{equation}

where $\bm{B}_1(\frac{d}{2})$ and $\bm{B}_2(\frac{d}{2})$ represent two butterfly matrices each of dimension $\frac{d}{2}$. Next, we define the \emph{butterfly component} as $\thickmuskip=2mu \medmuskip=2mu \tilde{\bm{B}}(d,k)=\text{diag}(\bm{B}^F_1(k),\cdots,\bm{B}^F_{d/k}(k))$, a block-diagonal matrix of size $d \times d$ with blocks of size $k$, where $\bm{B}^F_i(k), \forall i$ are butterfly factors as specified in Equation~\ref{eq:but_factor}. We are now prepared to utilize the butterfly matrix to parameterize an orthogonal matrix. To do this, it suffices to ensure that all multiplicative factors $\tilde{\bm{B}}(d,k)$, for every $k$, in the butterfly matrix $\bm{B}(d)$ are orthogonal. We first examine the block-diagonal matrix $\tilde{\bm{B}}(d,2)$ with block size $2$, and it is straightforward to enforce orthogonality on $\tilde{\bm{B}}(d,2)$ by parameterizing each block using the Cayley transform (or a 2-dimensional rotation), following the approach of \cite{liu2021orthogonal,qiu2023controlling}. Each butterfly component's non-zero structure can be interpreted as a permutation of the non-zero pattern of $\tilde{\bm{B}}(d,2)$, which means all butterfly components can be easily parameterized as orthogonal matrices. This approach provides an efficient parameterization of orthogonal matrices composed of numerous $2 \times 2$ orthogonal blocks. We extend the butterfly matrices following \cite{chen2022pixelated} by defining a block butterfly component $\tilde{\bm{B}}^b(d,k)$, where each scalar entry $\bm{d}_i$ is replaced by a $b \times b$ matrix. To ensure that the block butterfly component $\tilde{\bm{B}}^b(d,2)$ remains orthogonal, we parameterize each $2b \times 2b$ block matrix to be orthogonal. The non-zero patterns of other butterfly components $\tilde{\bm{B}}^b(d,k)$ for $k>2$ are block-wise permutations of the pattern in $\tilde{\bm{B}}^b(d,2)$, so they can similarly be converted into orthogonal matrices. Putting these elements together, the forward pass in BOFT is

\begin{equation}
    \begin{aligned}\nonumber
    \bm{z}=\big(\bm{R}(m,b)\cdot\bm{W}^0\big)^\top\bm{x},~~~~\text{s.t.}~\bigg{\{}\bm{R}(m,b)=\prod_{i=1}^m \tilde{\bm{B}}^b_{(d,i)}~~\&~~\underbrace{\big(\tilde{\bm{B}}^b_{(d,j)}\big)^\top\tilde{\bm{B}}^b_{(d,j)}=\tilde{\bm{B}}^b_{(d,j)}\big(\tilde{\bm{B}}^b_{(d,j)}\big)^\top\!=\bm{I}_d}_{\forall j\in [1, m]}\bigg{\}},
    \end{aligned}
\end{equation}

Here, for simplicity, we denote $\tilde{\bm{B}}^b(d,2^{m - i+1})$ as $\tilde{\bm{B}}^b_{(d,i)}$, and $\bm{I}_d$ represents the identity matrix of dimension $d$. The orthogonal matrix $\bm{R}(m,b)\in\mathbb{R}^{d\times d}$ is formed by multiplying several orthogonal butterfly components. For notational ease, BOFT with $\bm{R}(m,\frac{b}{2})$ is referred to as BOFT($m$,$b$), where $b\geq 2$. In the special case when $m=1$, BOFT($1$,$b$) simplifies to the block-diagonal OFT~\cite{qiu2023controlling} with block size $b$. Similarly, BOFT($1$,$d$) corresponds to the original OFT~\cite{qiu2023controlling} featuring an unconstrained full orthogonal matrix. Moreover, BOFT($\log \frac{2d}{b}$,$b$) employs the block butterfly matrix $\bm{B}^{\frac{b}{2}}(d)$ as $\bm{R}$, resulting in a dense orthogonal matrix $\bm{R}$. Generally, BOFT($m$,$b$) requires $\frac{1}{2}(b-1)dm$ trainable parameters to fine-tune a linear layer of dimensions $d\times n$. When the butterfly matrix is used, i.e., $m=\log d$ and $b=2$, BOFT consumes $\mathcal{O}(d\log d)$ parameters. In comparison, the original OFT with a dense full orthogonal matrix uses $\mathcal{O}(d^2)$ parameters, while the block-diagonal OFT with $r$ blocks uses $\mathcal{O}(bd)$. Therefore, to produce a dense orthogonal matrix, the original OFT must set $b=d$, whereas BOFT allows any $b$ to achieve this.

\textbf{Identity initialization for BOFT}. Typically, finetuning techniques commence from the exact pretrained model to ensure that the fine-tuned model remains close to the original. For instance, LoRA initializes its low-rank weights to zero. In BOFT, we initialize all butterfly components to identity matrices (that is, the skew-symmetric matrix within the Cayley transform is initialized to zero).

\textbf{Multiplicative Dropout for BOFT}. LoRA~\cite{hulora2022} applies a Dropout layer to the low-rank weight updates to mitigate overfitting. While standard Dropout~\cite{srivastava2014dropout} functions effectively for LoRA, it is incompatible with BOFT due to its multiplicative weight update scheme. To resolve this, we introduce a multiplicative Dropout tailored for BOFT. Since the orthogonal matrix $\bm{R}(m,b)$ comprises $m$ orthogonal butterfly components, which can be rearranged into $2b\times 2b$ block-diagonal orthogonal matrices, this multiplicative Dropout randomly selects $p_1$ percent of the butterfly components and $p_2$ percent of diagonal blocks within each component, substituting these selected blocks with identity matrices.

\section{Intriguing Insights and Discussions}

\textbf{Expressivity of BOFT}. The butterfly framework combined with permutations can exactly replicate several well-known fast linear transforms~\cite{dao2019learning,dao2019kaleidoscope} (such as the fast Fourier transform and Hadamard transform), yet the extent to which our orthogonal butterfly matrix can approximate a general orthogonal matrix is still unclear. To investigate this, we perform a simulation approximating a randomly generated dense orthogonal matrix~\cite{anderson1987generation} of size $512\times 512$. The outcomes presented in Figure~\ref{fig:para_eff} are averaged over 10 distinct random seeds. The y-axis indicates the approximation error, while the x-axis represents the count of effective trainable parameters. Each curve of the same color corresponds to a BOFT with a particular block size, where the leftmost point shows the error for BOFT($1$,$b$) (i.e., the original block-diagonal OFT with block size $b$). Generally, BOFT demonstrates superior parameter efficiency compared to OFT. For instance, BOFT($9$,$2$) outperforms BOFT($1$,$16$) in expressiveness while utilizing considerably fewer parameters. BOFT configurations with smaller $b$ and larger $m$ tend to be more parameter-efficient; for example, BOFT($6$,$4$) employs significantly fewer parameters but achieves a comparable approximation error to BOFT($2$,$16$). Overall, the butterfly matrix forms a more structured subset within the orthogonal group (relative to the block-diagonal structure), which theoretically makes BOFT strictly more expressive than OFT for the same block size.

\begin{theorem}[Expressivity of BOFT]\label{thm:exp_boft}
    BOFT possesses greater expressiveness than OFT with an equivalent block size. To approximate any orthogonal matrix of size $d$, one can compose butterfly matrices as $\bm{B}_{d-1,1}(d)\bm{B}^\top_{d-1,2}(d)\cdots\bm{B}_{1,1}(d)\bm{B}^\top_{1,2}(d)$, where each $\bm{B}_{i,j}(d)$ is a butterfly matrix for all indices $i,j$.
\end{theorem}

Theorem~\ref{thm:exp_boft} implies a straightforward extension of BOFT — the resulting orthogonal matrix can be generalized to $\thickmuskip=2mu \medmuskip=2mu \bm{R}^G(m_1,b_1,m_2,b_2,l)=\bm{R}_{l,1}(m_1,b_1)\bm{R}_{l,2}^T(m_2,b_2)\cdots\bm{R}_{1,1}(m_1,b_1)\bm{R}_{1,2}^T(m_2,b_2)$, where $\bm{R}_{i,j}^T(m,b)$ denotes the orthogonal matrix utilized in BOFT. When $m_1 = m_2 = \log d$, $b_1 = b_2 = 2$, and $l = d-1$, the matrix $\bm{R}^G(m_1,b_1,m_2,b_2,l)$ can represent the entire orthogonal group. This type of matrix construction is also referred to as a kaleidoscope hierarchy~\cite{dao2019kaleidoscope}. Nevertheless, it is important to note that enhanced expressiveness does not necessarily translate into improved finetuning performance, since full finetuning, despite its universal expressivity, often results in unsatisfactory outcomes. The balance between expressivity and regularity is crucial for the generalization capability during model finetuning. BOFT expands the finetuning parameter space with structural priors, thereby allowing a better balance between expressivity and regularity to be achieved.

\textbf{Spectral properties}. Orthogonal finetuning typically provides superior spectral characteristics compared to LoRA, as it exactly maintains the spectral norm of the pretrained weight matrix $\bm{W}^0$. This can be understood through singular value decomposition: $\bm{W}^0=\bm{U}\bm{\Sigma}\bm{V}^\top$, where $\bm{U},\bm{V}$ are orthogonal matrices and $\bm{\Sigma}$ is a diagonal matrix of singular values. Both OFT and BOFT apply an orthogonal matrix $\bm{R}$ on the left side, resulting in finetuned weights $\bm{R}\bm{U}\bm{\Sigma}\bm{V}^\top$, which does not alter the largest singular value (i.e., the spectral norm of $\bm{W}^0$). This preservation has been demonstrated to significantly enhance training stability and generalization~\cite{miyato2018spectral,yoshida2017spectral}. Additional interesting mathematical properties are provided in Appendix~\ref{app:sn_property}.

\textbf{Orthogonal finetuning as learning bilinear similarity}. BOFT can be expressed as learning the bilinear similarity $\bm{w}^0_i\bm{R}\bm{x}$, where $\bm{w}^0_i$ denotes the $i$-th neuron (i.e., column vector) of the weight matrix $\bm{W}^0$. BOFT can thus be interpreted as learning a bilinear similarity matrix $\bm{R}$ under a strong constraint (i.e., $\bm{R}$ must be orthogonal), which fundamentally relates to distance metric learning~\cite{xing2002distance} and bilinear forms~\cite{roman2005advanced}.

\textbf{Inductive bias and generalization in BOFT}. As $\bm{R}(m,b)$ in BOFT generally corresponds to a structured subset of the orthogonal group, imposing a restriction on the hypothesis space, BOFT inherently induces an inductive bias. We contend that the structured inductive bias imposed by butterfly factorization benefits generalization, given its shared structured patterns with numerous classical linear transforms~\cite{dao2019kaleidoscope}, including discrete Fourier transform, discrete sine/cosine transform, and Hadamard transform. Furthermore, the sparse matrix factorization employed in BOFT might also introduce an implicit inductive bias~\cite{gunasekar2017implicit,li2018algorithmic,liu2019neural}.

\textbf{Comparison to butterfly-based sparse training}. Several prior studies~\cite{dao2019learning,dao2019kaleidoscope,chen2022pixelated,dao2022monarch} explore sparse training utilizing the butterfly parameterization. These works generally concentrate on directly reparameterizing weight matrices through the butterfly structure and training neural networks from the beginning. The work in \cite{dao2022monarch} addresses finetuning pretrained weights by initially projecting them onto a variant of butterfly matrices and subsequently optimizing these projected components for downstream tasks. In contrast, BOFT introduces a markedly different finetuning approach that modifies the weights using weight matrices shared across layers.

\input{rebuttal-results}

\section{Concluding Remarks and Limitations}

This paper introduces Orthogonal Butterfly (BOFT), a versatile and parameter-efficient finetuning technique for foundation models grounded in the butterfly architecture. The central idea behind achieving enhanced parameter efficiency lies in representing a dense orthogonal matrix as a product of multiple sparse orthogonal matrices. To facilitate discovering feasible matrix factorizations, we present a graphical information transmission framework. Within this framework, we identify the butterfly structure as effectively satisfying our criteria for sparse orthogonal matrix factorization. We empirically validate BOFT’s efficacy in finetuning large language models, substantial vision models, and text-to-image generative models. Furthermore, our experiments confirm the advantages of BOFT as a broadly applicable finetuning method.

Despite its demonstrated empirical success, BOFT is not without limitations. Because the ultimate orthogonal matrix in BOFT results from the product of several orthogonal matrices, its training runtime incurs a modest overhead compared to OFT. Enhancing the training efficiency of BOFT remains an open challenge. Fortunately, following the finetuning phase, the orthogonal matrices learned via BOFT can be directly merged into the pretrained model, incurring no extra inference latency. Additionally, whether the butterfly network represents the optimal structure for information transmission is still uncertain. Our information transmission framework opens avenues to draw inspiration from a distinct field—computer networking—where the efficiency of network topologies in transmitting information is extensively examined. We anticipate that even more efficient network architectures could be leveraged for constructing dense orthogonal matrices.

\end{document}